% ============================================
% LUC-02
% Gerado de molambudos.md — build_miolo.py
% ============================================
\fragtitle{luc}{L}{LUC}{02}{O Arquivo}
\markboth{LUC-02 --- O Arquivo}{}

\noindent \textit{A carta de 1979 para 2026}

\noindent \textbf{Arquivo Público Mineiro --- Belo Horizonte}
\textbf{Fevereiro de 2026}

\vspace{0.5em}\noindent\rule{\textwidth}{0.5pt}\vspace{0.5em}

\noindent O Arquivo Público Mineiro fica num prédio cinza no centro de Belo Horizonte. Lúcia conhecia o lugar --- já tinha passado horas ali para outras pesquisas. Mas nunca tinha visto a sala dos fundos.

\noindent O diretor, Dr. Ronaldo Torres, a esperava na porta. Era um homem magro, de óculos grossos e mãos inquietas. Ele a cumprimentou com um aperto rápido.

\noindent \textit{\textquotedblleft{}Dra. Lúcia. Obrigado por vir tão rápido.\textquotedblright{}}

\noindent \textit{\textquotedblleft{}O senhor disse que era urgente.\textquotedblright{}}

\noindent \textit{\textquotedblleft{}É.\textquotedblright{}\textit{ Ele olhou para os lados. }\textquotedblleft{}Vamos para minha sala.\textquotedblright{}}

\noindent A sala do diretor era pequena, cheia de pastas e livros. Na mesa, um pacote envolto em pano preto e, ao lado, uma chave enferrujada sobre um papel de seda.

\noindent \textit{\textquotedblleft{}Chegou semana passada\textquotedblright{}\textit{, disse Ronaldo. }\textquotedblleft{}Uma mulher. Não deixou nome. Disse apenas: 'Isto pertence ao arquivo. O Dr. Heitor Oliveira pediu para devolver.'\textquotedblright{}}

\noindent \textit{\textquotedblleft{}Heitor Oliveira está desaparecido desde 1979.\textquotedblright{}}

\noindent \textit{\textquotedblleft{}Eu sei.\textquotedblright{}}

\noindent \textit{\textquotedblleft{}Como ele poderia pedir para devolver um diário quase cinquenta anos depois?\textquotedblright{}}

\noindent Ronaldo não respondeu. Ele tirou os óculos e limpou as lentes.

\noindent \textit{\textquotedblleft{}A senhora quer ver o diário?\textquotedblright{}}

\noindent \textit{\textquotedblleft{}Sim.\textquotedblright{}}

\noindent \textit{\textquotedblleft{}Antes, preciso mostrar uma coisa.\textquotedblright{}}

\noindent Ele abriu uma pasta sobre a mesa. Dentro, havia fotografias. Antigas. Amareladas. A primeira mostrava um menino magro, de braços finos e olho direito amarelo. Ele estava em pé, na frente de um muro de tijolos, com um número pendurado no pescoço: \textit{1.260}.

\begin{figure}[H]
  \centering
  \includegraphics[width=0.55\textwidth]{foto_01_paciente1260_sepia.png}
  \caption*{\footnotesize\textbf{Fotografia 1} --- \textit{Joaquim Antônio Correia --- Paciente 1.260. 1917.}}
\end{figure}

\noindent \textit{\textquotedblleft{}Joaquim\textquotedblright{}}, disse Lúcia.

\noindent \textit{\textquotedblleft{}Sim. Esta é a única foto dele. Foi tirada na chegada ao Colônia, em 1917.\textquotedblright{}}

\noindent \textit{\textquotedblleft{}E as outras?\textquotedblright{}}

\noindent Ronaldo virou as fotos. A segunda mostrava um homem --- um médico jovem, de jaleco, com olheiras profundas e um sorriso forçado.

\begin{figure}[H]
  \centering
  \includegraphics[width=0.55\textwidth]{foto_02_dr_heitor_sepia.png}
  \caption*{\footnotesize\textbf{Fotografia 2} --- \textit{Dr. Heitor Oliveira. Arquivo do Colônia, 1979.}}
\end{figure}

\noindent \textit{\textquotedblleft{}Dr. Heitor Oliveira. Foto do crachá, 1979.\textquotedblright{}}

\noindent A terceira foto fez Lúcia parar.

\noindent Era uma foto de um homem mais velho. Muito mais velho. Cabelos brancos, pele enrugada, olhos opacos. Ele estava sentado num banco de praça, com um livro no colo.

\begin{figure}[H]
  \centering
  \includegraphics[width=0.55\textwidth]{foto_03_homem_livro_sepia.png}
  \caption*{\footnotesize\textbf{Fotografia 3} --- \textit{Homem não identificado. Encontrada dentro do diário.}}
\end{figure}

\noindent \textit{\textquotedblleft{}Quem é?\textquotedblright{}} perguntou Lúcia.

\noindent \textit{\textquotedblleft{}Não sei\textquotedblright{}\textit{, disse Ronaldo. }\textquotedblleft{}Esta foto estava dentro do diário. Junto com uma carta.\textquotedblright{}}

\noindent \textit{\textquotedblleft{}Carta?\textquotedblright{}}

\noindent \textit{\textquotedblleft{}Endereçada à senhora.\textquotedblright{}}

\noindent Ronaldo pegou um envelope pardo. Estava lacrado. Escrito à mão:

\noindent \textit{\textquotedblleft{}Para Dra. Lúcia Mendes --- Pessoal\textquotedblright{}}

\noindent Lúcia apontou para a chave enferrujada sobre a mesa. \textit{\textquotedblleft{}E a chave?\textquotedblright{}}

\noindent Ronaldo pegou a chave com cuidado. \textit{\textquotedblleft{}Veio junto. Não sei se é do diário, do arquivo ou de alguma cela do Colônia. Já tentamos em todas as portas do prédio. Não encaixa em nenhuma.\textquotedblright{}}

\noindent \textit{\textquotedblleft{}Por que a mandaram junto?\textquotedblright{}}

\noindent \textit{\textquotedblleft{}A mulher que entregou o pacote disse: 'O Dr. Oliveira pedia para avisar que a chave não abre portas. Ela abre a história.'\textquotedblright{}}

\noindent Lúcia tocou a chave. O metal estava frio --- mais frio do que deveria estar numa sala fechada.

\noindent \textit{\textquotedblleft{}Guarde com o diário\textquotedblright{}\textit{, ela disse. }\textquotedblleft{}Pode ser importante.\textquotedblright{}}

\noindent \textit{\textquotedblleft{}Como sabiam que eu ia pedir o arquivo?\textquotedblright{}}

\noindent \textit{\textquotedblleft{}Não sei.\textquotedblright{}}

\noindent Lúcia abriu o envelope. Dentro, um papel amarelado, com caligrafia trêmula:

\vspace{0.5em}\noindent\rule{\textwidth}{0.5pt}\vspace{0.5em}

\noindent \textit{Dra. Lúcia,}

\noindent \textit{Se está lendo esta carta, o ciclo chegou até a senhora.}

\noindent \textit{Meu nome é Heitor Oliveira. Estou escrevendo esta carta em 2026. Ou melhor, estou escrevendo em 1979 e pedindo para que alguém entregue em 2026.}

\noindent \textit{Não pergunte como. O diário ensina.}

\noindent \textit{Li o diário do paciente 1.260 em março de 1979. Desde então, não sou mais o mesmo.}

\noindent \textit{Passei 47 anos tentando entender o que aconteceu comigo. Tentando entender o que aconteceu com Joaquim.}

\noindent \textit{Não entendi.}

\noindent \textit{Só entendi uma coisa: o diário não é um registro. É um veículo. A criatura que Joaquim encontrou na vala não mora em corpos. Mora em histórias. Cada vez que alguém lê, ela passa para o leitor.}

\noindent \textit{A senhora leu este diário?}

\noindent \textit{Não. Ainda não.}

\noindent \textit{Mas vai ler.}

\noindent \textit{E quando ler, vai entender.}

\noindent \textit{Heitor Oliveira}
\textit{--- 1979 (enviado para 2026)}

\vspace{0.5em}\noindent\rule{\textwidth}{0.5pt}\vspace{0.5em}

\noindent Lúcia colocou a carta na mesa. A mão tremia.

\noindent \textit{\textquotedblleft{}Isto é uma brincadeira\textquotedblright{}\textit{, ela disse. }\textquotedblleft{}Alguém está manipulando o arquivo.\textquotedblright{}}

\noindent \textit{\textquotedblleft{}Talvez\textquotedblright{}\textit{, disse Ronaldo. }\textquotedblleft{}Mas tem mais uma coisa.\textquotedblright{}}

\noindent \textit{\textquotedblleft{}O quê?\textquotedblright{}}

\noindent \textit{O diretor apontou para o pacote de pano preto. O diário.}

\noindent \textit{\textquotedblleft{}A mulher que entregou disse outra frase. Ela disse: 'Avise a Dra. Lúcia que o paciente 1.260 mandou um abraço.'\textquotedblright{}}

\noindent Lúcia engoliu seco.

\noindent Joaquim Antônio Correia morreu em 1979.

\noindent Isso era fato.

\noindent Mas o diário estava ali. Quente. Esperando.

\noindent \textit{\textquotedblleft{}Vou levar.\textquotedblright{}}

\noindent \textit{\textquotedblleft{}Tem certeza?\textquotedblright{}}, perguntou Ronaldo.

\noindent \textit{\textquotedblleft{}Não.\textquotedblright{}}

\noindent Ela estendeu a mão e tocou o pano preto.

\noindent O pano estava quente.

\vspace{0.5em}
\noindent\textit{\textbf{↪ Rotas:} \rota{LUC-03} \rotasep \rota{DOC-01} \rotasep \rota{DOC-20}}